\section{Binary trees and Catalan numbers}
\todo{Lecture 5}
\begin{theorem}
Let $r\in \mathbf{R}$. Then the function $(1+x)^{r}$ has a Taylor expansion
$$
(1+x)^{r}=\sum_{k=0}^{\infty} \binom{ r}  {k } x^{k}
$$
for $x \in (-1,1)$.
\end{theorem}

\begin{question}
Compute $\binom{1/3}  {2 }$.
\end{question}
\begin{answer}
$\binom{ 1/3}  {2 }=-\frac{1}{9}$.
\end{answer}

\begin{definition}[Binary tree]
An inductive definition of a binary tree is given as follows:

A binary tree is either empty, or consists of one distinguished vertex called the root, plus an ordered pair of binary trees called the left subtree and the right subtree.
\end{definition}

Let $b_{n}$ denote the number of binary trees with $n$ vertices. Our goal is to find a formula for $b_{n}$.
\begin{example}
$$
b_{0}=1\quad b_{1}=1\quad b_{2}=2
$$
\end{example}

The inductive definition of a binary tree implies the following recursive formula for $b_{n}$:
$$
b_{n}=b_{0}b_{n-1}+b_{1}b_{n-2}+\dots +b_{n-1}b_{0}
$$

Consider the generating function

$$
b(x)=\sum_{k=0}^{\infty} b_{n}x^{n}
$$
The recursive formula 
$$
b_{n}=\sum_{k=0}^{n-1} b_{k}b_{n-k-1}
$$
implies 
\begin{align*}
b(x)b(x) & =\left( \sum_{n=0}^{\infty} b_{n}x^{n} \right)\left( \sum_{n=0}^{\infty} b_{n}x^{n} \right) \\
 & =\sum_{k=0}^{\infty} \left( \sum_{m=0}^{k} b_{m}b_{k-m} \right)x^{k} \\
 & =\sum_{k=0}^{\infty} b_{k+1}x^{k}  \\
 & =\frac{1}{x}\left( \sum_{k=1}^{\infty} b_{k}x^{k} \right)=\frac{1}{x}(b(x)-b_{0})
\end{align*}
Thus $b(x)$ satisfies the identity
$$
xb^{2}(x)-b(x)+1=0
$$
which has two solutions
$$
b_{1,2}(x)= \frac{1\pm \sqrt{ 1-4x }}{2x}
$$
of which we will consider
$$
\tilde{b}(x)=\frac{1-\sqrt{ 1-4x }}{2x}=\frac{2}{1+\sqrt{ 1-4x }}
$$
that has a Taylor expansion around $x=0$
$$
\tilde{b}(x)=\sum_{n=0}^{\infty} \tilde{b}_{n}x^{n}
$$
and as $\tilde{b}(x)$ satisfies the equation 
$$
x\tilde{b}(x)^{2}-\tilde{b}(x)+1=0\Rightarrow \tilde{b}(0)=b_{0}=1
$$
which gives the sequence $\tilde{b}_{n}$:
$$
\tilde{b}_{n}=\tilde{b}_{0}\tilde{b}_{n-1}+\tilde{b}_{1}\tilde{b}_{n-2}+\dots +\tilde{b}_{n-1}\tilde{b}_{0}
$$
We have
\begin{align*}
b_{0}=1 &   =1 =\tilde{b}_{0} \\
b_{1}=b_{0}^{2}  & =\tilde{b}_{0}^{2}=\tilde{b}_{2} \\
b_{2}=b_{0}b_{1}+b_{1}b_{0} & =b_{0}b_{1}+b_{1}b_{0} =\tilde{b}_{2}  \\
 & ~~\vdots \\
\end{align*}
Therefore $b_{n}=\tilde{b}_{n}$.

To compute $\tilde{b}_{n}$ we use the generalized binomial theorem, which implies
$$
\sqrt{ 1-4x }=\sum_{n=0}^{\infty} (-4)^{n}\binom{ 1/2}  {n } x^{n}
$$
Thus
$$
\tilde{b}(x)=\sum_{n=1}^{\infty} -\frac{1}{2}(-4)^{n}\binom{ 1/2}  {n } x^{n-1} 
$$
therefore
$$
b_{n}=-\frac{1}{2}(-4)^{n+1}\binom{ 1/2}  {n+1 }
$$.

\begin{definition}
The numbers $b_{n}$ are called Catalan numbers.
\end{definition}
\begin{exercise}
Show that
$$
b_{n}=\frac{1}{n+1}\binom{ 2n}  {n }
$$
\end{exercise}

\subsection{More Catalan numbers}

\begin{definition}[Cutting polygons into triangles]
A diagonal triangulation of a polygon is a partition of the polygon into triangles by non-intersecting diagonals.
\end{definition}
\begin{lemma}
Let $n\ge3$. The number of diagonal triangulations a convex $n$-gon is the Catalan number $b_{n-2}$.
\end{lemma}
\begin{definition}[Parentheses]
A regular bracket structure satisfies the following two conditions:
\begin{enumerate}
\item The number of left and right brackets is the same.
\item The number of left brackets in any starting segment is not less than then number of right brackets.
\end{enumerate}
\end{definition}
\begin{lemma}
The number of regular bracket structures with $n$ pairs of brackets is the Catalan number $b_{n}$.
\end{lemma}
\begin{definition}[Ups and downs]
A Dyck path is a continuous broken line consisting of vectors $(1,1)$ and $(1,-1)$ starting at the origin $A=(0,0)$ and ending at the $x$-axis and always staying in the upper half-plane.
\end{definition}
\begin{lemma}
The number of Dyck paths consisting of $2n$ steps is the Catalan number $b_{n}$.
\end{lemma}

\section{Fibonnaci numbers and the Golden Ratio}

\begin{definition}
The Fibonnaci sequence $\{ F_{n} \}_{n=0}^{\infty}$ is defined by the following recursive formula:
$$
F_{0}=0,\quad F_{1}=1,\quad F_{n+1}=F_{n}+F_{n-1},~~ n \ge 1
$$
\end{definition}

Another interpretation of Fibonnaci numbers: Let $S_{n}$ denote the number of ways one can climb $n$ stairs if allowed to jump 1 or 2 stairs at once. $S_{n}$ is the number of solutions to the equation $x_{1}+x_{2}+\dots +x_{k}=n$ where $x_{i}\in \{ 1, 2\}$ and $k$ is not fixed. Thus we observe that
$$
S_{1}=1,\quad S_{2}=2,\quad S_{n}=S_{n-1}+S_{n-2}
$$
and $S_{n}=F_{n+1}$.

\begin{definition}[Golden ratio]
The golden ration is the number
$$
\varphi=\frac{1+\sqrt{ 5 }}{2}\approx1.61083\dots
$$
\end{definition}
\begin{proposition}
$$
\lim_{ n \to \infty } \frac{F_{n+1}}{F_{n}}=\varphi
$$
\end{proposition}

\subsection{Explicit formula for Fibonnaci numbers}

Consider the generating function of the Fibonnaci sequence :
$$
F(x)=F_{0}+F_{1}x+\dots +F_{n}x^{n}+\dots
$$
then
$$
x^{2}F(x)=\sum_{n=2}^{\infty} F_{n-2}x^{n},\quad xF(x)=\sum_{n=1}^{\infty} F_{n-1}x^{n} 
$$
thus
\begin{align*}
F(x)-xF(x)-x^{2}F(x) & =\sum_{n=0}^{\infty} F_{n}x^{n}-\sum_{n=1}^{\infty} F_{n-1}x^{n}-\sum_{n=2}^{\infty} F_{n-2}x^{n} \\
  & =F_{0}+xF_{1}-xF_{0}+\sum_{n=2}^{\infty} (F_{n}-F_{n-1}-F_{n-2})x^{n} \\
 & =x
\end{align*}
therefore
$$
F(x)(1-x-x^{2})=x
$$
thus
$$
F(x)=\frac{x}{1-x-x^{2}}
$$
as a formal power series.

The function $\frac{x}{1-x-x^{2}}$ has derivatives of all orders at the point $x=0$. We will compute the Taylor expansion of this function at $x=0$. 

The polynomial $1-x-x^{2}$ has two roots
$$
x_{1,2}=\frac{-1\pm \sqrt{ 5 }}{2}
$$
therefore
\begin{align*}
 \frac{x}{1-x-^{2}} & =-\frac{x}{(x-x_{1})(x-x_{2})} \\
 &=\left( \frac{1}{x-x_{1}}-\frac{1}{x-x_{2}} \right) \frac{x}{x_{2}-x_{1}} \\
 & =\frac{x}{\sqrt{ 5 }}\left( -\frac{1}{x_{1}}\cdot \frac{1}{1-\frac{x}{x_{1}}}+\frac{1}{x_{2}}\cdot \frac{1}{1-\frac{x}{x_{2}}}  \right) \\
 & =\frac{x}{\sqrt{ 5 }}\left( \sum_{n=0}^{\infty}  -\frac{1}{x_{1}} \left( \frac{x}{x_{1}} \right)^{n}+\sum_{n=0}^{\infty}  \frac{1}{x_{2}}\left( \frac{x}{x_{2}} \right)^{n}\right) \\
 & =\frac{1}{\sqrt{ 5 }}\sum_{n=0}^{\infty} (-x_{1}^{-n-1}+x_{2}^{-n_-1})x^{n+1} \\
 & =\sum_{n=1}^{\infty}  \frac{1}{\sqrt{ 5 }}\left( \left( \frac{1+\sqrt{ 5 }}{2} \right)^{n}-\left( \frac{1-\sqrt{ 5 }}{2} \right)^{n} \right)x^{n}
\end{align*}
Therefore the Fibonnaci numbers satisfy
$$
F_{n}=\frac{1}{\sqrt{ 5 }}\left( \left( \frac{1+\sqrt{ 5 }}{2} \right)^{n}-\left( \frac{1-\sqrt{ 5 }}{2} \right)^{n} \right)=\frac{1}{\sqrt{ 5 }}(\varphi^{n}-(-\varphi)^{-n})
$$
for $n\ge 0$.